\chapter*{Résumé}
\markboth{Résumé}{}
\addcontentsline{toc}{chapter}{Résumé}
\adjustmtc
\thispagestyle{MyStyle}

Ce rapport de projet de fin d'études présente le développement d'une plateforme de dématérialisation de la gestion des candidatures, destinée à digitaliser l'organisation des concours : gestion des candidats, des centres, des établissements et des salles, répartition automatique des candidats, puis génération et envoi des convocations.

Ce projet a été réalisé dans le cadre de mon stage de fin d'études au sein du Ministère de l'Économie et des Finances, à Rabat. Le système vise à assister les gestionnaires dans le traitement des candidatures et, surtout, à automatiser l'affectation des candidats aux salles d'examen selon les capacités disponibles, tout en fiabilisant l'envoi des convocations et en offrant une vision opérationnelle via un tableau de bord.

Il apporte au ministère un outil intégré qui remplace un traitement encore largement manuel --- tableurs, affectations sujettes aux erreurs et envois fragmentés --- par une chaîne unique, de l'import des candidats jusqu'à la convocation. La valeur ajoutée réside dans le gain de temps pour les gestionnaires, la réduction des erreurs d'affectation (capacité des salles, double placement, éloignement géographique) et la fiabilisation de la notification des candidats, tout en conservant un contrôle humain. Pour y parvenir, la solution a été conçue puis développée sous forme d'application web interne : une interface destinée aux agents, un backend organisé en microservices Spring Boot et des bases PostgreSQL, avec une répartition automatique selon la capacité des salles et la proximité géographique, la génération de convocations PDF et un envoi groupé par courrier électronique.

Ainsi, ce résumé situe le projet dans son cadre de digitalisation de la gestion des concours et dans son contexte de réalisation au Ministère de l'Économie et des Finances, tout en soulignant sa portée opérationnelle pour les gestionnaires.

\begin{spacing}{1.5}
\underline{\textbf{Mots-clés :}} dématérialisation, gestion des concours, répartition automatique, convocations, microservices, Spring Boot, React, PostgreSQL, JWT.
\end{spacing}
